\subsection{10 августа. Дневка под пер. Иттиши}

\textit{Метеоусловия: ясно, во второй половине дня дождь, после сильный град}

При написании отчёта автор ориентировался на данные хронометриста. Вот эти данные. 

\begin{figure}[h!]
	\centering
	\includegraphics[angle=0, width=0.7\linewidth]{../pics/days/10_1.png}
	\label{fig:10_1}
	\caption{Отчёт хронометриста на 10 августа.}
\end{figure}


Подъём. Он был, точно не рано, т.к. группа была уставшая.  Когда -уже никогда доподлинно не будет известно. Днёвке предшествовал не очень сложный, но достаточно стрессовый подъём на перевал, который состоялся довольно поздно. 
Чем-то вкусным покушали, что-то вкусное выпили. Утро было максимально ненапряжным, т.к. погода была отличная, а первый контент видеограф записал в 13:09 по местному времени.
Когда-то был обед.
Много купались (рядом было озеро, холодное, но с рыбой), отдыхали и расслаблялись, восстанавливаясь после пройденного. У Дарьи Мерзликиной наконец-то начала проходить горная болезнь, что не могло не радовать.
В 17:45 на горизонте показалась группа Остапива Алексея, в связи с чем руковод и штурман прямо из озера в одних купательных принадлежностях побежали их встречать. Встреча прошла успешно, группа Лёши встала с палатками рядом с нами.
День не ознаменовался крупными подвигами или рекордами пройденного расстояния.
Впрочем, именно в этот день была сделана одна из лучших фотографий похода, выигравшая в последствии фотоконкурс в паблике CatTours за август 2025 г. – «Геометрическая гора».

\begin{figure}[h!]
	\centering
	\includegraphics[angle=0, width=0.7\linewidth]{../pics/days/10_2.jpg}
	\label{fig:10_2}
	\caption{<<Геометрическая гора>>}
\end{figure}

Примерно в 18.00 устроили спортивные мероприятия. Руковод и штурман стояли в планке, позднее к ним присоединился и Шалфеев Илья, простоявший в планке всю песню «Warriors of the world» группы Manowar.

\begin{figure}[h!]
	\centering
	\includegraphics[angle=0, width=0.7\linewidth]{../pics/days/10_3.png}
	\label{fig:10_3}
	\caption{Руковод и штурман стоят в планке, Хронометрист фиксирует время.}
\end{figure}


Когда-то был ужин, вкусный.
Когда-то был отбой. 


\begin{tabular}{|c|c|}
	\hline
	ЧХВ: & 0:00 \\
	\hline
	ГХВ: & 0:00 \\
	\hline
\end{tabular}
\clearpage